\section*{Tag 11 — Datamatrix-Anatomie und ASCII-Encoding}

\subsection*{Bleistift-Übung 1: Bauteile am HONIG-Symbol}

Beim 12$\times$12-Datamatrix mit Inhalt \texttt{HONIG} (siehe Buch
oben):

\begin{itemize}
\item Das \textbf{L-Pattern} besteht aus der \emph{kompletten unteren
  Kante} (12 schwarze Module nebeneinander) und der \emph{kompletten
  linken Kante} (12 schwarze Module übereinander). Zusammen ergeben
  sie das „L".
\item Das \textbf{Zeitsignal} liegt gegenüber: die obere Kante und
  die rechte Kante haben \emph{abwechselnd} schwarze und weiße
  Module — schwarz, weiß, schwarz, weiß, \dots
\item Der \textbf{Datenbereich} ist alles dazwischen: ein 10$\times$10-
  Innenrechteck, in dem die 12 Codewörter (5 Daten + 7 Prüfbytes)
  als 96 Module verteilt sind. Welches Modul wo landet, regelt
  der Datenplatzierungs-Algorithmus aus Etappe 13.
\end{itemize}

Eine kleine Faustregel: das L-Pattern wirkt wie ein „Sitzplatz" —
der Boden und die Lehne sind voll schwarz. Davor sitzt das Datenfeld,
und das Zeitsignal ist die abwechselnde Polster-Naht oben und rechts.

\subsection*{Bleistift-Übung 2: HONIG zu Codewort-Bytes}

Der ASCII-Wert plus 1 ergibt das Datamatrix-Codewort:

\begin{center}
\small
\begin{tabular}{c c c}
\toprule
Buchstabe & ASCII-Wert & Codewort (ASCII + 1) \\
\midrule
H & 72 & 73 \\
O & 79 & 80 \\
N & 78 & 79 \\
I & 73 & 74 \\
G & 71 & 72 \\
\bottomrule
\end{tabular}
\end{center}

Codewort-Folge für \texttt{HONIG}: \texttt{[73, 80, 79, 74, 72]}.

Das ist genau die Eingabe für den Reed-Solomon-Encoder, den wir in
Etappe 12 anwenden.

\subsection*{Bleistift-Übung 3: Greta vorbereiten}

Beispiel: das Wort \texttt{GRETA}.

\begin{center}
\small
\begin{tabular}{c c c}
\toprule
Buchstabe & ASCII-Wert & Codewort (ASCII + 1) \\
\midrule
G & 71 & 72 \\
R & 82 & 83 \\
E & 69 & 70 \\
T & 84 & 85 \\
A & 65 & 66 \\
\bottomrule
\end{tabular}
\end{center}

Codewort-Folge für \texttt{GRETA}: \texttt{[72, 83, 70, 85, 66]}.

Das ist deine Datenfolge für Etappe 12, wenn du das eigene Symbol
zeichnen willst.
